% !TeX root = ../教材题库_全部条目_转写工作稿.tex
% Source: XJ-B1-C04-LAB-P137-Y0145
% 第4章 / 4.4 函数与方程 / 教材 p.137、138 / PDF 第144、145页
% OCR未校对：公式与文字均可能有误；未完成数学TeX转写；禁止直接发给学生
% 原题及相邻原解截图（核对用）：../题目原图/XJ-B1-C04-LAB-P137-Y0145-01.webp, ../题目原图/XJ-B1-C04-LAB-P137-Y0145-02.webp
% 本文件只在转写完成后作为题干使用；OCR 条目尚不是可靠题干。
\begin{bankitem}{用二分法求方程的近似解}{第4章 / 4.4 函数与方程 / 教材 p.137、138 / PDF 第144、145页}{XJ-B1-C04-LAB-P137-Y0145}{OCR未校对：公式与文字均可能有误；未完成数学TeX转写；禁止直接发给学生}
\begin{ocrdraft}
[教材p.137 / PDF第144页]
\par
数学实验
\par
用二分法求方程的近似解
\par
实际问题：如图1（a），有一个轴线水平放置的直径为2m的圆柱形储油罐，
\par
当
\par
储油量为油罐容积的一时，油的深度是多少？
\par
0
\par
(a)
\par
问题可以化为更一般的数学问题：如图1（b），弦，AB把半径为R的圆分成两
\par
块，要使较小一块的面积是圆面积的
\par
，如何确定弦AB到圆心的距离？
\par
设AB的中点为M,ZAOM=k0，则弦AB到圆心的距离OM一Rcos差。．这时，
\par
差丌R
\par
扇形OAB的面积是
\par
，而AOAB的面积是。，于是，问题化为求满
\par
180
\par
足下列等式的数差：
\par
差丌R
\par
—Rsink0cos差
\par
180
\par
化简后得到关于未知数差的方程：
\par
a差丌
\par
F（差，a)==vr十asink0cos差
\par
180
\par
下面根据实际问题取磊一3，记不差）一F（差，3），得到方程
\par
丆（D一丌十3sink0cos旷
\par
60
\par
设方程的解为，以下分三步进行：
\par
第一步，确定差、的大致范围．取差一60，得丆（60）一299，“>0，取差一90，得
\par
不90）一一l.57小．<0，可见e（60，90）．
\par
第二步，通过二分法迭代计算可得出e（74，6，74．7).暂将区间（74．6，
\par
74，7）的中点作为解的近似值，即取74．65\ensuremath{\cdot}
\par
第三步，在原问题中，=1m，当ZAOM一74．65。时，有OM=Rcos74．65。
\par
[教材p.138 / PDF第145页]
\par
0．2647m．
\par
又知道的准确值在74．6到74，7之间．由于Rcos74．6。0．2656m，
\par
Rcos74．7。0．2639m,可见我们的解答结果的实际误差在1mm之内．
\par
因此，
\par
油的深度为1十0．2647=1.2647（m)．
\par
使用“网络画板”或“超级画板”中的表达式测量和动画按钮赋值功能，
\par
可以
\par
用鼠标操作控制二分法计算步骤，如图2\ensuremath{\cdot}
\par
用二分法求函数零点近似值丆（0一丌+3sin々。co。一缸一0
\par
60
\par
[：@丝云钽1]到b一预设误差界限0彐
\par
@0和些鲎0！」」《加
\par
口=．53/2b亠73750区/长之2。巧\&
\par
衤在《三@《《@《@申在《]=m0770567
\par
9次，屮，点一
\par
／4巧30丿=\#59巧9
\par
2
\par
若一0则m即为所求零点，结束．
\par
当区间长度之半小于预设误差界限时，中点m即为所求零点的所要近似值，结束．
\par
如图3，左边是二分法的通用程序，右边是它在“超级画板”的程序工作区运
\par
行4．4，2小节例4中的问题的情形．
\par
F10at(1);
\par
F10at(l);
\par
>>计算结果显示洱点数\#
\par
f（对\{\ensuremath{\times}\ensuremath{\cdot}3一3蔫\ensuremath{\cdot}2；》
\par
f仆；\}
\par
0h11e（u〈0\&\&b—a>d)
\par
\{c以ab）／2；VFf(C);
\par
while（u〈0b-a〉d)
\par
else\{（u区\{u竺"；b-c；\}
\par
else\{u-w;a-c
\par
if伟一的\{u：喏》
\par
else\{if（u〈0）\{v=w,
\par
>>以对
\par
else\{u=w;a=c;\}\}\}
\par
（a期/2；\}
\par
eff（0，1，0．001）；
\par
>>（133/）／（208）=0．652832囍
\end{ocrdraft}
\end{bankitem}
